% ============================================================================
% CHAPTER 2 TEACHER'S MANUAL
% Complete instructional guide for teaching AI uncertainty concepts
% ============================================================================
% Source: /markdown/ch2/Ch2T_updated.md
% Note: This file is for the Instructor's Edition, not the main book
% ============================================================================

\chapter{Teacher's Manual: Chapter 2}
\label{ch:teacher-ch02}

\begin{chapterquote}{Complete instructional guide}
Teaching AI uncertainty concepts using \emph{When Smart Systems Get Confused}.
\end{chapterquote}

% ============================================================================
\section{Course Integration Guide}
% ============================================================================

\subsection{Learning Objectives}

By the end of this chapter, students will be able to:

\paragraph{Knowledge (Remembering \& Understanding).}
\begin{itemize}
    \item Define aleatoric and epistemic uncertainty and give examples of each
    \item Explain what ``calibration'' means in the context of AI confidence scores
    \item Describe why a system can be wrong without being confused, and confused without being wrong
    \item Identify three ways systems can signal (or fail to signal) their uncertainty
\end{itemize}

\paragraph{Application \& Analysis.}
\begin{itemize}
    \item Classify real-world AI errors as aleatoric vs.\ epistemic in origin
    \item Evaluate whether a system's uncertainty signals are actionable or merely noisy
    \item Apply the ``confusion-is-useful'' framing to analyze HITL design choices
    \item Assess how training data gaps manifest as predictable confusion patterns
\end{itemize}

\paragraph{Synthesis \& Evaluation.}
\begin{itemize}
    \item Design uncertainty signaling strategies for different HITL applications
    \item Critique AI systems for calibration failures and propose improvements
    \item Argue for or against adding human review based on the type of uncertainty present
    \item Evaluate the HITL value of human review given aleatoric vs.\ epistemic uncertainty sources
\end{itemize}

\subsection{Prerequisites}
\begin{itemize}
    \item \textbf{Chapter~1 concepts:} The Five Dimensions framework, the three failure modes, alert fatigue
    \item \textbf{Basic probability:} Concept of probability distributions
    \item \textbf{Not required:} Machine learning background (main chapter); Appendix~2A requires basic statistics
\end{itemize}

\subsection{Course Positioning}

\begin{description}
    \item[Early course (Week 2--3)] Best used immediately after Chapter~1 --- builds the conceptual vocabulary for the rest of the book
    \item[Standalone] Can be taught as ``Types of AI Failure'' in a broader AI literacy course
\end{description}

% ============================================================================
\section{Key Concepts for Instruction}
% ============================================================================

\subsection{The Central Reframe}

The chapter's central argument is a reframe: AI ``confusion'' is not a failure to avoid but a \emph{signal to exploit}. The pedagogical goal is to shift students from ``a good AI never gets confused'' to ``a good AI knows when it's confused.''

\begin{technote}[Teaching Tip]
Open with Dr.~Anand's story before introducing any vocabulary. Ask students what the AI did right \emph{before} telling them. Most students will initially say the AI failed (it didn't diagnose correctly). The reveal --- that the 51\% score was the valuable information --- creates a memorable cognitive shift.
\end{technote}

\subsection{Aleatoric vs.\ Epistemic Uncertainty}

\begin{table}[htbp]
\caption{Aleatoric vs.\ epistemic uncertainty: core comparison}
\label{tab:uncertainty-types}
\centering
\begin{tabular}{@{}p{3cm}p{4cm}p{4cm}@{}}
\toprule
\textbf{Property} & \textbf{Aleatoric} & \textbf{Epistemic} \\
\midrule
Source & Noise in the world & Model's ignorance \\
Reducible? & No & Yes (more data) \\
Human review helps? & Not reliably & Usually yes \\
Example & Blurry photograph & Unknown species \\
HITL response & Acknowledge ambiguity & Route to expert \\
\bottomrule
\end{tabular}
\end{table}

\begin{keyconcept}[Critical Distinction]
Students often conflate ``the model is uncertain'' with ``a human would also be uncertain.'' The key distinction is: aleatoric uncertainty = \emph{anyone} would be uncertain; epistemic uncertainty = \emph{this model} is uncertain, but an expert might not be. Only epistemic uncertainty reliably benefits from human review.
\end{keyconcept}

\subsection{The Calibration Concept}

Calibration is the chapter's most technical concept and requires care in presentation. Recommended approach:

\begin{enumerate}
    \item Start with the intuitive version: ``If your friend says `I'm 90\% sure' about things, they should be right 9 times out of 10. If they're right only 6 times out of 10, their confidence is miscalibrated.''
    \item Apply to AI: the same principle, but measurable quantitatively.
    \item Connect to HITL: poor calibration means the routing signal is unreliable.
\end{enumerate}

\begin{cautionbox}[Key Misconception to Preempt]
Students often assume a well-calibrated model is a highly accurate model. Clarify: a model can be 60\% accurate and perfectly calibrated (it always says ``I'm 60\% sure''), or 90\% accurate and poorly calibrated (it always says ``I'm 99\% sure''). Calibration is about the \emph{relationship} between confidence and accuracy, not accuracy alone.
\end{cautionbox}

\subsection{The 2×2 Wrong/Confused Matrix}

Use this framework to organize the four possible system states:

\begin{table}[htbp]
\caption{The wrong/confused matrix: four system states}
\label{tab:wrong-confused}
\centering
\begin{tabular}{@{}lll@{}}
\toprule
 & \textbf{High Confidence} & \textbf{Low Confidence} \\
\midrule
\textbf{Correct} & Best case (trust it) & Useful signal (review) \\
\textbf{Wrong} & Most dangerous (silent error) & Good signal (routes review) \\
\bottomrule
\end{tabular}
\end{table}

The most dangerous quadrant is confident + wrong: there is no signal to catch the error. The most \emph{valuable} quadrant for HITL purposes is uncertain (low confidence), whether right or wrong, because it triggers appropriate review.

% ============================================================================
\section{Lecture Planning Guide}
% ============================================================================

\subsection{50-Minute Lecture Structure}

\subsubsection{Opening Story: The Radiologist (8 minutes)}

Tell Dr.\ Anand's story without revealing the punchline (the 51\% score) until you've asked the class what happened.

\begin{trythis}
\textbf{Poll:} ``Did the AI succeed or fail at its job?''\\[0.5em]
\textbf{Reveal:} The AI succeeded at its \emph{actual} job --- expressing calibrated uncertainty --- even though it didn't produce a diagnosis.\\[0.5em]
\textbf{Transition:} ``This chapter is about understanding what AI confusion actually looks like and why it's valuable.''
\end{trythis}

\subsubsection{Part 1: Two Kinds of Uncertainty (12 minutes)}

\begin{itemize}
    \item Introduce aleatoric vs.\ epistemic with the frosted-glass-photo/Patagonian-mara contrast
    \item \textbf{Class activity (3 min):} Give three examples; have students vote aleatoric vs.\ epistemic:
    \begin{enumerate}
        \item A speech recognizer failing on a noisy recording (mostly aleatoric)
        \item A spam filter struggling with a new type of phishing email (mostly epistemic)
        \item A face recognizer failing on dark-skinned faces (epistemic --- training data gap)
    \end{enumerate}
    \item \textbf{Key insight:} Only epistemic uncertainty reliably benefits from human review
\end{itemize}

\subsubsection{Part 2: The Calibration Problem (10 minutes)}

\begin{itemize}
    \item Introduce calibration with the ``trusted friend'' analogy
    \item Show reliability diagram concept: plot confidence vs.\ accuracy; perfect calibration is the diagonal
    \item \textbf{Example:} A medical AI that says ``90\% confident'' but is only right 70\% of the time
    \item \textbf{HITL connection:} If confidence scores are miscalibrated, routing cases to human review happens on the wrong cases
\end{itemize}

\subsubsection{Part 3: How Confusion Surfaces (or Doesn't) (10 minutes)}

Three modes: explicit signals, behavioral signals, silent failures.

\begin{table}[htbp]
\caption{Comparison of uncertainty surfacing modes}
\label{tab:surfacing-modes}
\centering
\begin{tabular}{@{}lll@{}}
\toprule
\textbf{System} & \textbf{Mode} & \textbf{Outcome} \\
\midrule
Medical AI (Ch.\ 2) & Explicit confidence score & Caught rare infection \\
Air Canada chatbot (Ch.\ 1) & Silent failure & Wrong answer with full confidence \\
Voice recognition & Silent failure & Errors without uncertainty flags \\
\bottomrule
\end{tabular}
\end{table}

\begin{keyconcept}[Core Principle]
A system that fails silently breaks the HITL loop before it can begin. The value of an uncertainty flag equals its precision: too many flags creates alert fatigue; too few creates overconfidence.
\end{keyconcept}

\subsubsection{Part 4: Wrong vs.\ Confused (5 minutes)}

\begin{itemize}
    \item Present the $2\times2$ matrix (Table~\ref{tab:wrong-confused})
    \item Most dangerous: confident + wrong (no signal to catch)
    \item Most valuable signal: uncertain (routes to review appropriately in all cases)
\end{itemize}

\subsubsection{Wrap-Up and Preview (5 minutes)}

\begin{itemize}
    \item Confusion is a feature, not a bug --- if it's surfaced
    \item \textbf{Assign ``Try This''} exercise: monitor confidence signals in daily technology use
    \item \textbf{Preview:} ``Next chapter: where does this confusion come from? Understanding how AI systems learn explains why they get confused in predictable patterns.''
\end{itemize}

\subsection{75-Minute Extended Session}

\paragraph{Add: Group Activity --- Uncertainty Type Audit (15 min).}
Small groups analyze one AI application (spam filter, medical imaging, voice assistant, fraud detection, content moderation). Task: identify likely sources of uncertainty (aleatoric vs.\ epistemic), assess how the system signals confusion, propose one calibration improvement.

\paragraph{Add: Technical Interlude (10 min).}
Walk through reliability diagram construction with a simple numeric example. Show what overconfidence looks like in a diagram. Connect ECE (from Appendix~2A) to practical implications.

% ============================================================================
\section{Discussion Questions by Level}
% ============================================================================

\subsection{Introductory Level}

\begin{enumerate}
    \item \textbf{Recognition:} ``Dr.~Anand's story shows an AI that didn't give the right diagnosis but still saved the patient. How? What exactly did the AI do correctly?''\\
    \textit{Target insight: It expressed calibrated uncertainty, which triggered human review.}

    \item \textbf{Concept Check:} ``Explain the difference between aleatoric and epistemic uncertainty in your own words. Give one example of each from your own experience with technology.''\\
    \textit{Target insight: Aleatoric = noise in the world; epistemic = model's knowledge gap.}

    \item \textbf{Comparison:} ``Compare how the medical imaging AI in Chapter~2 and the Air Canada chatbot in Chapter~1 handled their respective uncertainties. Which worked better? Why?''\\
    \textit{Target insight: Medical AI surfaced uncertainty explicitly; chatbot buried it → silent failure.}

    \item \textbf{Application:} ``A voice recognition system has 40\% word error rate on Scottish accents but only 5\% on Midwestern American accents. What type of uncertainty is this? Should the system flag its output differently for Scottish speakers? How?''\\
    \textit{Target insight: Epistemic (training data gap). Yes --- selective flagging for underrepresented speech patterns.}
\end{enumerate}

\subsection{Intermediate Level}

\begin{enumerate}
    \item \textbf{Framework Application:} ``Apply the Five Dimensions framework to the radiology AI scenario. Which dimensions are well-implemented? Which might be weak?''\\
    \textit{Guide: Uncertainty Detection (excellent --- expressed 51\%); Stakes Calibration (presumably good --- medical setting); Feedback Integration (unclear --- does the radiologist's correction improve the model?).}

    \item \textbf{Design Thinking:} ``Design an uncertainty signaling strategy for a legal document review AI. The AI reads contracts and flags potential issues. When should it express uncertainty? How?''\\
    \textit{Target: Distinguish routine clauses (high confidence, silent) from unusual language (low confidence, flagged with specific reason).}

    \item \textbf{Calibration Analysis:} ``A fraud detection system reports 95\% confidence on 10,000 transactions. It's wrong on 800 of them. Is this system well-calibrated? What does this mean for HITL routing?''\\
    \textit{Answer: Expected accuracy at 95\% confidence = 95\%, actual = 92\%. ECE \approx 0.03 for this bin. Reasonably calibrated, but systematic --- worth temperature scaling.}

    \item \textbf{Trade-off Analysis:} ``An image classifier for content moderation is 99\% accurate on 85\% of images, but drops to 65\% accurate on the remaining 15\%. Should the HITL system route all 15\% to human review, or try to distinguish within that group?''\\
    \textit{Target: Within the 15\%, further stratification by confidence would route the most uncertain subset. Blanket routing wastes human capacity.}
\end{enumerate}

\subsection{Advanced Level}

\begin{enumerate}
    \item \textbf{Technical Design:} ``Using the uncertainty decomposition from Appendix~2A, how would you determine whether a medical imaging AI's errors are primarily aleatoric or epistemic? What would each finding imply for HITL design?''

    \item \textbf{Calibration Strategy:} ``A model trained on one hospital's patient population is deployed at a different hospital. How would you expect its calibration to change? What calibration correction methods would you apply, and on what data?''

    \item \textbf{System Design:} ``Design a complete uncertainty pipeline for a spam filter, from confidence estimation through HITL routing. Specify: uncertainty estimation method, calibration check, routing threshold, interface for human reviewer, and feedback integration.''

    \item \textbf{Research Extension:} ``The chapter argues that `confusion reveals training data gaps.' Design a methodology for using uncertainty patterns to prioritize which data to collect next for model improvement.''
\end{enumerate}

% ============================================================================
\section{Hands-On Activities}
% ============================================================================

\subsection{Activity 1: Uncertainty Type Sorting (10 minutes)}

\textbf{Setup:} Cards or slides with AI failure scenarios.

\textbf{Task:} Students sort into aleatoric, epistemic, or mixed.

\paragraph{Scenarios.}
\begin{description}
    \item[Storm recording] Speech recognizer failing on a recording in a storm (aleatoric)
    \item[New product category] Image classifier failing on a product category never trained on (epistemic)
    \item[Sarcasm] Sentiment analyzer failing on sarcasm (epistemic --- pattern not learned)
    \item[Two conditions] Medical AI failing on a patient with two simultaneous rare conditions (mixed)
    \item[Untranslatable] Translation system failing on a word with no direct equivalent (epistemic)
    \item[Legitimate-unusual] Fraud detector uncertain on a transaction that is simultaneously legitimate and unusual (aleatoric)
\end{description}

\textbf{Debrief:} For each case, ``Would a human expert be uncertain here too?'' If yes → aleatoric component. If no → epistemic.

\subsection{Activity 2: Calibration Intuition (15 minutes)}

\textbf{Setup:} Historical examples with real accuracy rates and stated confidence levels.

\begin{table}[htbp]
\caption{Calibration comparison across systems}
\label{tab:calibration-comparison}
\centering
\begin{tabular}{@{}p{4cm}p{3cm}p{3cm}p{3cm}@{}}
\toprule
\textbf{System} & \textbf{Stated Confidence} & \textbf{Actual Accuracy} & \textbf{Assessment} \\
\midrule
Weather app ``90\% chance of rain'' & 90\% & 85\% & Slightly overconfident \\
Spam filter ``high confidence'' & 95\% & 95\% & Well calibrated \\
Medical chatbot ``I'm sure this is\ldots'' & 99\% & 72\% & Severely overconfident \\
Navigation ``arrives in 15 min'' & $\sim$95\% & 70\% & Overconfident \\
\bottomrule
\end{tabular}
\end{table}

\textbf{Task:} Students calculate the calibration gap and rank from best to worst calibrated. Discuss which system creates the most dangerous HITL failure.

\subsection{Activity 3: Design Challenge --- Surfacing Confusion (20 minutes)}

\textbf{Setup:} Small groups.

\textbf{Scenario:} A college uses an AI system to read admissions essays and give each a score from 1--10. The AI is 85\% accurate when confident, but accuracy drops to 62\% when the essay is non-standard (highly creative, non-native English, unusual background).

\paragraph{Tasks.}
\begin{enumerate}
    \item What type of uncertainty applies to the 15\% ``hard'' cases?
    \item Design an uncertainty signaling system for this AI
    \item Specify: what threshold triggers human review? What information does the reviewer see?
    \item How do you prevent alert fatigue (assuming reviewers can handle $\sim$20\% of essays)?
\end{enumerate}

\textbf{Deliverable:} One-page design spec + justification.

\subsection{Activity 4: Five-Dimension Framework Audit (20 minutes)}

\textbf{Task:} Apply the Five Dimensions framework to the radiology AI scenario from Chapter~2's opening.

\begin{table}[htbp]
\caption{Five-dimension framework worksheet (radiology AI)}
\label{tab:radiology-audit}
\centering
\begin{tabular}{@{}p{3.5cm}p{2cm}p{3.5cm}p{3.5cm}@{}}
\toprule
\textbf{Dimension} & \textbf{Rating (1--5)} & \textbf{Evidence} & \textbf{Improvement} \\
\midrule
Uncertainty Detection & & & \\[8pt]
Intervention Design & & & \\[8pt]
Timing & & & \\[8pt]
Stakes Calibration & & & \\[8pt]
Feedback Integration & & & \\
\bottomrule
\end{tabular}
\end{table}

\textbf{Discussion:} Which dimension is most commonly ignored in real deployed systems?

% ============================================================================
\section{Common Student Misconceptions}
% ============================================================================

\subsection{Misconception 1: ``Better AI Means No Uncertainty''}

\paragraph{Student thinking:} ``If AI is good enough, it won't be confused.''

\paragraph{Correction strategy.}
\begin{itemize}
    \item Aleatoric uncertainty is irreducible --- no AI can be certain about a blurry photograph
    \item For HITL design, some uncertainty is \emph{desired} --- it's the routing mechanism
    \item The goal is calibrated uncertainty, not zero uncertainty
\end{itemize}

\subsection{Misconception 2: ``Low Confidence = Wrong Answer''}

\paragraph{Student thinking:} ``If the AI isn't sure, I should ignore it.''

\paragraph{Correction strategy.}
\begin{itemize}
    \item A 51\% confidence score means ``I'm slightly more likely right than wrong'' --- not ``I'm wrong''
    \item The value is in routing: low-confidence outputs go to human review, improving overall accuracy
    \item Dr.~Anand's AI was uncertain about a case that turned out to be genuinely difficult
\end{itemize}

\subsection{Misconception 3: ``Calibration Equals Accuracy''}

\paragraph{Student thinking:} ``A 90\%-accurate model is well-calibrated.''

\paragraph{Correction strategy.}
\begin{itemize}
    \item Use the weather forecaster analogy: someone who always says ``50\% chance of rain'' can be right 50\% of the time (calibrated) or 100\% of the time in a rainy climate (miscalibrated-underconfident)
    \item Calibration is about the \emph{relationship} between expressed confidence and accuracy
\end{itemize}

\subsection{Misconception 4: ``Human Review Always Helps''}

\paragraph{Student thinking:} ``When in doubt, ask a human.''

\paragraph{Correction strategy.}
\begin{itemize}
    \item For aleatoric uncertainty, human review doesn't reliably improve outcomes --- the ambiguity is in the world
    \item Human review costs time and attention; over-routing creates alert fatigue
    \item Chapter~4 formalizes this: human review should be routed to where human knowledge adds value (epistemic cases) and error costs are high
\end{itemize}

\subsection{Misconception 5: ``Silent Systems Are More Trustworthy''}

\paragraph{Student thinking:} ``A confident-looking AI seems more reliable.''

\paragraph{Correction strategy.}
\begin{itemize}
    \item The Air Canada chatbot \emph{looked} confident and was wrong
    \item Systems that express uncertainty honestly are more trustworthy, not less
    \item Research shows users who understand AI uncertainty make better decisions than users who assume infallibility
\end{itemize}

% ============================================================================
\section{Assessment Options}
% ============================================================================

\subsection{Option 1: Uncertainty Audit Paper (500--750 words)}

\textbf{Prompt:} Choose an AI-powered product you use. Analyze:
\begin{enumerate}
    \item What types of uncertainty does it likely encounter (aleatoric vs.\ epistemic)?
    \item How does it currently signal (or fail to signal) confusion?
    \item What would better uncertainty signaling look like?
    \item How would improved signaling enable better HITL design?
\end{enumerate}

\paragraph{Rubric.}
\begin{itemize}
    \item \textbf{Uncertainty type analysis (30\%):} Correct application of aleatoric/epistemic distinction
    \item \textbf{Signal assessment (30\%):} Specific, accurate description of current behavior
    \item \textbf{Design proposal (25\%):} Feasible, well-reasoned improvement
    \item \textbf{Framework connection (15\%):} Links to Five Dimensions
\end{itemize}

\subsection{Option 2: Calibration Case Study}

Given provided confidence score and accuracy data for a real or simulated system:
\begin{enumerate}
    \item Construct a reliability diagram
    \item Calculate ECE
    \item Recommend a calibration correction method
    \item Describe the HITL routing implications of the calibration gap
\end{enumerate}

\subsection{Option 3: Discussion Post + Peer Response}

\textbf{Post:} 200-word response to: ``A medical AI reports 95\% confidence on a diagnosis but recent tests show it's only 80\% accurate on similar cases. Should you trust it? What should happen next?''

\textbf{Peer response:} Engage with a classmate's post, either building on or challenging their analysis.

% ============================================================================
\section{Extension Activities}
% ============================================================================

\subsection{Individual Research Topics}
\begin{enumerate}
    \item \textbf{Historical Analysis:} How have weather forecasting systems improved calibration over the last 50 years? What lessons apply to AI systems?
    \item \textbf{Domain Study:} Research calibration practices in financial risk modeling. How do they compare to typical ML deployments?
    \item \textbf{Cross-Cultural:} How do different cultures interpret probabilistic forecasts? Does this affect how uncertainty signals should be designed?
\end{enumerate}

\subsection{Group Projects}
\begin{enumerate}
    \item \textbf{Calibration Audit:} As a group, test the calibration of a publicly available AI service. Design the test, collect data, analyze results.
    \item \textbf{Case Database:} Collect documented cases of AI systems that failed silently (overconfident wrong outputs). Analyze each for which type of uncertainty was present and whether HITL would have helped.
\end{enumerate}

% ============================================================================
\section{Connections to Other Chapters}
% ============================================================================

\subsection{Backward Links}
\begin{itemize}
    \item \textbf{Chapter~1:} Alert fatigue (too many flags = calibration problem); the Goldilocks problem; uncertainty quantification introduced
\end{itemize}

\subsection{Forward Links}
\begin{itemize}
    \item \textbf{Chapter~3:} Where does uncertainty come from? Training data gaps, overfitting, the train/test gap
    \item \textbf{Chapter~4:} Threshold problem --- what to do with a confidence score; false positives vs.\ false negatives
    \item \textbf{Chapter~5:} Interface design for uncertainty communication --- how to present the flag to the human
\end{itemize}

\subsection{Cross-Curricular Connections}
\begin{itemize}
    \item \textbf{Statistics:} Probability calibration; reliability diagrams; expected calibration error
    \item \textbf{Philosophy of Science:} Epistemic vs.\ aleatory uncertainty in scientific reasoning
    \item \textbf{Medical Informatics:} Diagnostic uncertainty communication; clinical decision support
    \item \textbf{Psychology:} How humans interpret probabilistic information; overconfidence in human judgment
\end{itemize}

% ============================================================================
\section{Sample Lesson Plan}
% ============================================================================

\begin{table}[htbp]
\caption{Sample 50-minute lesson plan for Chapter~2}
\label{tab:lesson-plan-ch02}
\centering
\begin{tabular}{@{}p{2.5cm}p{6cm}p{4cm}@{}}
\toprule
\textbf{Time} & \textbf{Activity} & \textbf{Materials} \\
\midrule
0:00--0:08 & Opening story: radiologist + poll & Ch.\ 2 slide with X-ray scenario \\
0:08--0:20 & Aleatoric vs.\ epistemic + sorting activity & Scenario cards or slide \\
0:20--0:30 & Calibration concept + reliability diagram & Diagram slide + numeric example \\
0:30--0:40 & Surfacing confusion: 3 modes + comparisons & Case study slides \\
0:40--0:45 & Wrong vs.\ confused 2×2 & Whiteboard or slide \\
0:45--0:50 & Wrap-up, assign Try This, preview Ch.\ 3 & Chapter handout \\
\bottomrule
\end{tabular}
\end{table}

\bigskip
\begin{center}
\rule{0.5\textwidth}{0.4pt}
\end{center}
\bigskip

\noindent\textit{This teacher's manual provides everything needed to successfully teach Chapter~2 on AI uncertainty. The central reframe --- confusion as signal rather than failure --- is the conceptual foundation students need for the threshold, interface, and annotation chapters that follow.}

\medskip
\noindent\textit{Next: Chapter~3 explores where AI uncertainty originates --- in the training process itself --- and why that makes HITL demand predictable.}
